\section{The four classes}\label{sec:classes}

\subsection{Conventions}

Categories are small unless they are categories of sets, groups or modules.
For a category $A$ we write $NA$ for its nerve and $BA$ for its classifying
space. A functor is a \emph{weak equivalence} if it induces a weak homotopy
equivalence of classifying spaces, and $A$ is \emph{weakly contractible} if
$BA$ is. A natural transformation between two functors induces a homotopy
between the induced maps, and a category with an initial object is
contractible \cite[\S 1, Proposition~2 and Corollary~2]{Quillen}. For a
functor $u\colon A\to B$ and an object $b$ of $B$, the objects of
$b\backslash u$ are the pairs $(a,\ b\to u(a))$. The functor $u$ is
\emph{final} if every $b\backslash u$ is non-empty and connected
\cite[IX.3]{CWM}. Quillen's Theorem~A \cite[\S 1]{Quillen}: if every
$b\backslash u$ is weakly contractible, then $u$ is a weak equivalence.

For a functor $L\colon A\to\Set$ the \emph{category of elements} $\El(L)$ has
as objects the pairs $(a,l)$ with $l\in L(a)$, and as arrows
$(a,l)\to(a',l')$ the arrows $f\colon a\to a'$ with $L(f)l=l'$. An
\emph{action component} of $A$ is a connected component of $\El(L)$ for some
$L$. If $D$ is an action component of $A$, a component of $\El(L)$, and
$K\colon D\to\Set$ is a functor, then
\begin{equation}\label{eq:composite}
 \El(K)\cong\El(\widetilde K),\qquad
 \widetilde K(a)=\coprod_{l\in L(a),\ (a,l)\in D}K(a,l) ;
\end{equation}
so every action component of $D$ is an action component of $A$. If
$\Psi\colon D\to\Set$ sends all arrows to bijections, then
$B\El(\Psi)\to BD$ is a covering space, and every covering space of $BD$
arises in this way \cite[\S 1, Proposition~1]{Quillen}.

For a functor $N\colon A\to\Ab$ let $C_\bullet(A;N)$ be the chain complex with
$C_n(A;N)=\bigoplus_{a_0\to\dots\to a_n}N(a_0)$, the sum being over the
$n$-simplices of $NA$, and with the usual alternating differential; let
$H_n(A;N)$ be its homology. We use three properties: $H_0(A;N)=\colim_AN$; a
short exact sequence of functors gives a long exact sequence of the groups
$H_n$; and for $L\colon A\to\Set$ the complex $C_\bullet(A;\Z[L])$ is the
complex of integral chains of $N\El(L)$. In particular, for
$P_x=\Z[A(x,-)]$,
\begin{equation}\label{eq:rep}
 H_0(A;P_x)=\Z,\qquad H_n(A;P_x)=0\quad(n>0),
\end{equation}
because $\El(A(x,-))=x\backslash A$ has an initial object; the isomorphism
$H_0(A;P_x)\cong\Z$ is induced by the augmentation
$\sum n_if_i\mapsto\sum n_i$.

\subsection{Definitions}

\begin{definition}\label{def:classes}
Let $A$ be a non-empty connected category.
\begin{enumerate}
\item[\textup{(a)}] $A$ is \emph{filtered} if any two objects have a common
 successor and for every pair $f,g\colon p\to q$ there is an arrow $h$ with
 $hf=hg$.
\item[\textup{(b)}] $A$ is \emph{exact} if the functor $\colim_A\colon\Ab^A\to\Ab$ is
 exact.
\item[\textup{(c)}] $A$ is \emph{total} if for every final functor $u\colon A\to B$
 the category $B$ is weakly contractible.
\item[\textup{(d)}] $A$ is \emph{quasi-filtered} if \textup{(QF1)} every pair of
 arrows $q\leftarrow p\to q'$ can be completed to a commutative square, and
 \textup{(QF2)} for every pair $f,g\colon p\to q$ there are an object $r$ and
 an element $t$ of $\Z[A(q,r)]$ of augmentation $1$ with $tf=tg$ in
 $\Z[A(p,r)]$.
\end{enumerate}
\end{definition}

Exact categories were characterized by Isbell and Mitchell
\cite{IM73,IM76}; an element $t$ as in (QF2) is an \emph{affine equalizer} in
their terminology, and conditions of the form (QF1), (QF2) are studied in
\cite{IM76} and, under the name quasi-filtered, by Ad\'amek, Koubek and
Reiterman \cite{AKR}. We use these two conditions only as stated here.
Property (c) is Grothendieck's \emph{asph\'ericit\'e totale}, for the weak
homotopy equivalences \cite[Chapter~II, ms.~p.~36ter; Chapter~XII,
D\'efinition~1]{Der}. It is not the property called totally aspherical in
the theory of test categories and in \cite{Rosicky}, which concerns the
diagonal $A\to A\times A$; see Proposition~\ref{prop:sifted}.

\subsection{Total categories}

\begin{lemma}\label{lem:acyclic}
Let $A$ be a category all of whose action components have the integral
homology of a point. Then all action components of $A$ are weakly
contractible.
\end{lemma}

\begin{proof}
Let $D$ be an action component and $\gamma$ an element of the fundamental
group of $BD$. The connected covering space of $BD$ whose fundamental group
is the cyclic group generated by $\gamma$ is $B\El(\Psi)$ for some
$\Psi\colon D\to\Set$, and $\El(\Psi)$ is an action component of $A$ by
\eqref{eq:composite}. Its first homology group is that cyclic group, so
$\gamma=1$. Thus $BD$ is simply connected with the homology of a point, and
it is weakly contractible by the Hurewicz theorem.
\end{proof}

\begin{proposition}\label{prop:total}
A non-empty connected category $A$ is total if and only if every action
component of $A$ is weakly contractible.
\end{proposition}

\begin{proof}
Suppose that all action components are weakly contractible, and let
$u\colon A\to B$ be final. For an object $b$, the category $b\backslash u$ is
the category of elements of the functor $a\mapsto B(b,u(a))$, and it is
connected; so it is an action component and is weakly contractible. By
Theorem~A, $u$ is a weak equivalence. Since $A$ is the category of elements
of the functor with value a point, $BA$ is weakly contractible, and so is
$BB$.

Conversely let $A$ be total, and let $D$ be an action component, a component
of $\El(L)$. Let $L_D\subseteq L$ be the subfunctor with
$L_D(a)=\{l:(a,l)\in D\}$. Let $A_D$ be the category obtained from $A$ by
adjoining one object $v$, with $A_D(v,a)=L_D(a)$, $A_D(a,v)=\varnothing$,
$A_D(v,v)=\{\id\}$, and with composition given by the action. The inclusion
$A\to A_D$ is final: for an object $a$ of $A$ the category $a\backslash A$
has an initial object, and $v\backslash A=D$. Hence $A_D$ is weakly
contractible, and so is $A$, because the identity of $A$ is final. A simplex
of $NA_D$ that is not in $NA$ is a chain that begins with $v$; so $NA_D$ is
the union of $NA$ and of the nerve of the category $D^{\triangleleft}$
obtained from $D$ by adjoining an initial object, glued along $ND$ by the
projection $ND\to NA$. Hence $BA_D$ is the mapping cone of $BD\to BA$, and
the reduced integral homology of $BD$ vanishes. By Lemma~\ref{lem:acyclic}
all action components are weakly contractible.
\end{proof}

\begin{remark}\label{rem:hocolim}
Equivalently, $A$ is total if and only if for every functor $X$ from $A$ to
simplicial sets the canonical map $\operatorname{hocolim}_AX\to\colim_AX$ is a
weak equivalence. Indeed, the homotopy colimit is the diagonal of the
bisimplicial set $(m,n)\mapsto N_n\El(X_m)$ \cite[XII.5.2]{BK}, and the map
is induced by the maps $N\El(X_m)\to\pi_0\El(X_m)=\colim_AX_m$. If the action
components are weakly contractible these are weak equivalences, and then so is
the map of diagonals \cite[IV, Proposition~1.7]{GJ}. Conversely a functor
$L\colon A\to\Set$ is a diagram of discrete simplicial sets, for which the map
is $N\El(L)\to\pi_0\El(L)$.
\end{remark}

\subsection{The inclusions}

\begin{proposition}\label{prop:incl}
Let $A$ be a non-empty connected category. If $A$ is filtered it is exact,
and if $A$ is exact it is total.
\end{proposition}

\begin{proof}
Filtered colimits of abelian groups are exact
\cite[Theorem~2.6.15]{Weibel}. Let $A$ be exact, and let $D$ be an action
component, a component of $\El(L)$. For $N\colon D\to\Ab$ put
$(\Sigma N)(a)=\bigoplus_{l:(a,l)\in D}N(a,l)$. The functor
$\Sigma\colon\Ab^D\to\Ab^A$ is exact and $\colim_A\Sigma N=\colim_DN$, so
$\colim_D$ is exact. Let $R_\bullet$ be the complex of functors $D\to\Ab$
with $R_n=\bigoplus_{d_0\to\dots\to d_n}\Z[D(d_n,-)]$ whose value at an object
$d$ is the complex of integral chains of the nerve of $D/d$. As $D/d$ has a
terminal object, $R_\bullet$ is a resolution of the constant functor $\Z$.
Its colimit over $D$ is the complex of integral chains of $ND$. Since
$\colim_D$ is exact, $BD$ has the integral homology of a point. By
Lemma~\ref{lem:acyclic} and Proposition~\ref{prop:total}, $A$ is total.
\end{proof}

\begin{lemma}\label{lem:succ}
Let $D$ be a connected category such that $H_1(B\El(K);\Z)=0$ for every
functor $K\colon D\to\Set$ with $\El(K)$ connected. Then any two objects of
$D$ have a common successor.
\end{lemma}

\begin{proof}
Let $x,x'$ be objects, and let $U$ and $U'$ be the full subcategories of the
objects that receive an arrow from $x$, respectively from $x'$. They are
cosieves. Suppose that they are disjoint and put $V=U\cup U'$. Let
$K(d)=\{1,2\}$ for $d\notin V$ and $K(d)=\{1\}$ for $d\in V$, with the
evident maps. Then $N\El(K)$ is the union of two copies of $ND$ that meet in
$NV$, and $\El(K)$ is connected. In the Mayer--Vietoris sequence
\[
 H_1(B\El(K);\Z)\longrightarrow H_0(BV;\Z)\longrightarrow
 H_0(BD;\Z)\oplus H_0(BD;\Z)
\]
the middle group has rank at least $2$, since no arrow joins $U$ and $U'$,
and the image of the second map has rank $1$. So the first group is not
zero, a contradiction.
\end{proof}

\begin{proposition}\label{prop:qf}
Let $A$ be a non-empty connected category such that $H_1(BD;\Z)=0$ for every
action component $D$ of $A$. Then $A$ is quasi-filtered. In particular every
total category is quasi-filtered.
\end{proposition}

\begin{proof}
By \eqref{eq:composite}, Lemma~\ref{lem:succ} applies to every action
component of $A$.

(QF1). Let $f\colon p\to q$ and $f'\colon p\to q'$. The category
$p\backslash A$ is an action component. A common successor of its objects
$(q,f)$ and $(q',f')$ is a commutative square $uf=u'f'$.

(QF2). Let $f,g\colon p\to q$. Let $H$ be the quotient of the functor
$A(p,-)$ by the equivalence relation generated by $kf\sim kg$ for all arrows
$k$ with source $q$. Then $\El(H)$ is connected, because every object
receives an arrow from $(p,[\id_p])$. There is an exact sequence
\[
 P_q\xrightarrow{\ \partial\ }P_p\longrightarrow\Z[H]\longrightarrow 0,
 \qquad\partial(k)=kf-kg .
\]
Let $K=\ker\partial$ and $J=\operatorname{im}\partial$. The sequences
$0\to J\to P_p\to\Z[H]\to0$ and $0\to K\to P_q\to J\to 0$ and \eqref{eq:rep}
give
\[
 H_1(A;\Z[H])=\ker\bigl(\colim J\to\colim P_p\bigr),\qquad
 \colim J=\operatorname{coker}\bigl(\colim K\to\colim P_q=\Z\bigr).
\]
The elements of $J$ have augmentation zero, so $\colim J\to\colim P_p=\Z$ is
zero. Hence $H_1(B\El(H);\Z)=H_1(A;\Z[H])\cong\Z/I$, where $I$ is the subgroup
of $\Z$ generated by the augmentations of the elements of the groups
$K(x)=\{t\in\Z[A(q,x)]:tf=tg\}$. By hypothesis $I=\Z$: there are
$t_j\in K(x_j)$ and integers $n_j$ with $\sum_jn_j\,\varepsilon(t_j)=1$. By
Lemma~\ref{lem:succ}, applied to $A$, there are arrows $u_j\colon x_j\to r$
into one object $r$, and $t=\sum_jn_ju_jt_j$ has augmentation $1$ and
satisfies $tf=tg$.

The last assertion follows from Proposition~\ref{prop:total}.
\end{proof}

\subsection{Sifted categories}

A category $A$ is \emph{sifted} if it is non-empty and the diagonal
$\Delta\colon A\to A\times A$ is final \cite{AR}.

\begin{proposition}\label{prop:sifted}
A non-empty connected category is filtered if and only if it is total and
sifted.
\end{proposition}

\begin{proof}
Let $A$ be total and sifted. For objects $a,b$ the category
$(a,b)\backslash\Delta$ is non-empty, so $a$ and $b$ have a common successor.
Let $f,g\colon p\to q$. The category $(p,p)\backslash\Delta$ is the category
of elements of $A(p,-)\times A(p,-)$, and it is connected; so it is an action
component of $A$. By Proposition~\ref{prop:total} and Lemma~\ref{lem:succ}
its objects $(q,f,g)$ and $(q,f,f)$ have a common successor $(r,w,w')$:
there are $h,h'\colon q\to r$ with $hf=w=h'f$ and $hg=w'=h'f$. Hence $hf=hg$,
and $A$ is filtered.

Conversely a filtered category is total by Proposition~\ref{prop:incl}, and
it is sifted: two objects $(c,u,v)$ and $(c',u',v')$ of
$(a,b)\backslash\Delta$ are joined through an object $(d,\cdot,\cdot)$, where
$c\to d\leftarrow c'$ are chosen so as to equalize the two resulting arrows
$a\to d$ and the two resulting arrows $b\to d$.
\end{proof}
